\section{Consequences}\label{sec:free-by-cyclic}

Let $G$ be a hyperbolic one-relator group or a hyperbolic ascending
mapping torus of an injective endomorphism of a finite-rank free group.
Theorems~\ref{thm:main} and~\ref{thm:ascending-intro} give a
compact-special finite-index subgroup $H\leq G$. The conclusions below
also apply to a finitely generated Baumslag--Solitar-free one-relator
group: the companion hyperbolicity theorem
\cite[Theorem~1.1]{OpenAIHyperbolicity} first puts it in the first class. We explain the finite-quotient and subgroup consequences before
the more specific virtual-fibering result.

\subsection{Linearity and finite quotients}

The finite Salvetti model embeds $H$ in a right-angled Artin group on a
finite graph. Such groups are linear over $\mathbb Z$
\cite{HaglundWise2008}. Inducing a faithful integral representation
from the finite-index subgroup $H$ gives a faithful representation of
$G$ in $\mathrm{GL}_N(\mathbb Z)$ for some finite $N$. Reduction of
matrix entries modulo suitable integers then distinguishes every
nonidentity element in a finite quotient. Thus $G$ is linear over
$\mathbb Z$ and residually finite. The conclusion concerns the whole
group, including the one-relator groups with torsion; it is not merely
a property of the chosen cover. Borisov--Sapir had already proved
residual finiteness for ascending free-group tori without requiring
hyperbolicity \cite{BorisovSapir2005}; the compact-special result here
also supplies integral linearity.
For finite-rank free bases, this answers Dru\c{t}u--Sapir's question
about the existence of hyperbolic nonlinear ascending HNN extensions
in the negative \cite[Problem~6]{DrutuSapir2005}.

Finite quotients also distinguish conjugacy classes. A group is
\emph{conjugacy separable} if any two nonconjugate elements have
nonconjugate images in some finite quotient, and \emph{hereditarily
conjugacy separable} if every finite-index subgroup has this property.
Minasyan--Zalesskii's theorem applies to every hyperbolic virtually
compact special group, so $G$ is hereditarily conjugacy separable
\cite[Theorem~1.1]{MinasyanZalesskii2016}. For one-relator groups with
torsion, they had already established this conclusion
\cite[Theorem~1.1]{MinasyanZalesskii2013}.

There is also a cohomological consequence of the finite special model.
Canonical completion realizes $H$ as a virtual retract of a
right-angled Artin group: it is a retract of a finite-index subgroup of
that group. Proposition~3.8 of
\cite{MinasyanZalesskii2016} therefore makes $G$ \emph{cohomologically
good}: for every finite $G$-module $M$ and every $n\geq0$, the natural
comparison map
\[
H^n_{\mathrm{cont}}(\widehat G;M)\longrightarrow H^n(G;M)
\]
is an isomorphism. Here $\widehat G$ is the profinite completion and the
left side uses continuous cochains. Thus the finite-quotient structure
recovers group cohomology with finite coefficients, not just the
ability to distinguish individual elements.

\subsection{Quasiconvex subgroups and largeness}

If $Q\leq G$ is quasiconvex, then $Q_0=Q\cap H$ is quasiconvex in $H$.
Haglund--Wise's canonical completion and retraction give a finite-index
subgroup $H_0\leq H$ containing $Q_0$ and a retraction
$H_0\to Q_0$ \cite[Section~6 and Theorem~7.3]{HaglundWise2008}.
In particular, $Q_0$ is separable in $H$: every element of
$H\setminus Q_0$ can be excluded by a finite-index subgroup containing
$Q_0$. Subgroup separability passes through this finite-index change,
so $Q$ is separable in $G$ \cite[Lemma~7.5]{HaglundWise2008}.
This permits finite covers that retain a specified quasiconvex
subgroup while excluding any prescribed finite set outside it. The
retraction gives the stronger structural control on $Q\cap H$ stated
above. The quasiconvexity hypothesis is essential to this application;
finite generation alone is not being substituted for it.

When $G$ is \emph{non-elementary}, meaning that it is not virtually
cyclic, it is also \emph{large}: some finite-index subgroup surjects
onto a free group of rank two. Indeed, the torsion-free hyperbolic
subgroup $H$ contains a quasiconvex free subgroup of rank two; one may
apply \cite[Theorem~1.11]{KielakLinton2024Virtually} to $H$ with the
trivial subgroup as the prescribed family. The virtual retraction just
described then supplies the surjection. In particular, finite-index
subgroups of $G$ have arbitrarily large rational first Betti number: pull back
finite-index subgroups of arbitrarily large rank in that free quotient.
The non-elementarity condition is an additional hypothesis for these
largeness and homology conclusions.

\subsection{Virtual fibering and ambient embeddings}

Agol's virtual RFRS theorem for right-angled Artin groups, together
with the Salvetti embedding and inheritance by subgroups, gives a
finite-index RFRS subgroup of $G$ \cite[Corollary~2.3]{Agol2008}. In this condition, a descending sequence
$H'=H'_0\geq H'_1\geq\cdots$ of finite-index subgroups, normal in
$H'$, has trivial intersection, and each $H'_{i+1}$ contains the kernel
of $H'_i\to H_1(H'_i;\mathbb Q)$. It connects finite covers with
rational homology and underlies the virtual-fibering results discussed
in the introduction. For the ascending tori,
Corollary~\ref{cor:ascending-finite-rank} combines it with field
acyclicity to give a finite-index epimorphism onto $\mathbb Z$ with
finite-rank free kernel. General one-relator groups need not have that
acyclicity. The theorems of Kielak--Linton instead give the following
subgroup-closed conclusion and finite-rank ambient model. The
one-relator conclusion is the implication they explicitly anticipated
for Wise's Conjecture~17.8 \cite[Section~1.1]{KielakLinton2024Virtually}.

\begin{corollary}\label{cor:free-by-cyclic}
Let $G$ be one of the following groups.
\begin{enumerate}
\item A one-relator group
\[
G=F(X)/\langle\!\langle r\rangle\!\rangle,
\]
where $X$ is finite and $r\in F(X)$ is arbitrary, such that either $G$
is word-hyperbolic or $G$ contains no subgroup isomorphic to
$\mathrm{BS}(m,n)$ for any nonzero integers $m,n$.
\item A word-hyperbolic ascending mapping torus
\[
G=\langle F,t\mid t^{-1}ft=\phi(f)\ (f\in F)\rangle,
\]
where $F$ is a free group of finite rank and $\phi:F\to F$ is injective.
\end{enumerate}
There are a finite-index subgroup $H\leq G$, an integer $d\geq0$, an
automorphism $\alpha$ of the rank-$d$ free group $F_d$, and an embedding
\[
H\hookrightarrow F_d\rtimes_\alpha\mathbb Z.
\]
In particular, $G$ is virtually free-by-cyclic: the same $H$ fits into
an exact sequence
\[
1\longrightarrow K\longrightarrow H\longrightarrow Q\longrightarrow1,
\qquad K\text{ free},\quad Q\leq\mathbb Z.
\]
The rank of $K$ may be infinite, and $Q$ may be trivial.
\end{corollary}

\begin{proof}
If $G$ is finite, take $H=1$ and $d=0$, so the ambient group is
$\mathbb Z$. We may therefore assume that $G$ is infinite.

For the first case, a compact connected planar surface with $|X|+1$
boundary components has fundamental group $F(X)$; for $X=\varnothing$
this surface is a disk. If $G$ is Baumslag--Solitar-free, Theorem~1.1
of the companion \cite{OpenAIHyperbolicity} first makes it hyperbolic.
Theorem~\ref{thm:main} then gives virtual compact specialness.
Kielak--Linton's Corollary~1.3 \cite{KielakLinton2024Virtually}
applies to the quotient of this surface
group by the normal closure of $r$, and makes $G$ virtually
free-by-cyclic. Its relator is arbitrary, so this application does not
exclude trivial relators or proper powers.

In the second case, Corollary~\ref{cor:ascending-finite-rank} already
gives a finite-index free-by-$\mathbb Z$ subgroup with finite-rank
free kernel. The weaker virtual-free-by-cyclic conclusion also
follows from Theorem~\ref{thm:ascending-intro} and Kielak--Linton's
Corollary~1.8 \cite{KielakLinton2024Virtually}.

For the stronger embedding, apply Theorem~1.1 of
\cite{KielakLinton2024Virtually} to the hyperbolic virtually compact
special group $G$. Its equivalence of (b) and (c) gives a finite-index
subgroup embedded in a free-by-cyclic group with finitely generated
free kernel.
The proof of that theorem gives an ambient quotient $\mathbb Z$ in its
rational cohomological dimension two case. In its dimension-at-most-one
case it instead supplies a finite-index free subgroup; that subgroup has
finite rank because $G$ is finitely generated, and it embeds in its
direct product with $\mathbb Z$. Thus in either case the ambient group
has an exact sequence with finite-rank free kernel and quotient
$\mathbb Z$. Choosing a lift of $1\in\mathbb Z$ splits the sequence and
gives the displayed $F_d\rtimes_\alpha\mathbb Z$.

Identify $H$ with its image in this semidirect product and let
$\pi:F_d\rtimes_\alpha\mathbb Z\to\mathbb Z$ be the projection. Its
restriction to $H$ has kernel $K=H\cap F_d$ and image $Q=\pi(H)$.
The subgroup theorem for free groups makes $K$ free, and $Q$ is a
subgroup of $\mathbb Z$. This proves the final exact sequence. Only the
ambient kernel $F_d$ is asserted to be finitely generated; finite
generation does not follow for its subgroup $K$.
\end{proof}

For an ascending torus, one can use the stronger finite-index subgroup
from Corollary~\ref{cor:ascending-finite-rank} as its own ambient
finite-rank free-by-$\mathbb Z$ group. Its fibering character is not
identified with the canonical ascending height. For injective
nonsurjective ascending tori, this answers the
hyperbolic case of Kielak's virtual-fibering question recorded by
Mutanguha \cite[Overview, p.~3, Question (Virtual fibering)]{Mutanguha2021Dynamics}.
The separate quadratic-Dehn-function variant in that source is not
covered by the hyperbolicity hypothesis here. Thus a finite-index
subgroup is a mapping torus of a finite-rank free-group automorphism,
even when the original ascending map is not surjective. This stronger
conclusion is specific to the acyclic ascending setting; the general
one-relator conclusion retains the possibly infinite-rank kernel and
possibly trivial quotient stated above.
